%---------------------------------------------------------------------
% Title page and one orientation page.
%---------------------------------------------------------------------
\thispagestyle{empty}

% Band depth measured from the rendered page, with headroom. Nastaliq at
% Scale=1.45 sets far taller lines than the nominal point sizes suggest.
\begin{tikzpicture}[remember picture, overlay]
  \fill[bmblued] (current page.north west) rectangle
        ([yshift=-12.3cm]current page.north east);
  \fill[bmaccent] ([yshift=-12.3cm]current page.north west) rectangle
        ([yshift=-12.5cm]current page.north east);
\end{tikzpicture}

\vspace*{0.9cm}
{\color{white}
  {\large علامہ اقبال اوپن یونیورسٹی، اسلام آباد}\par
  \vspace{3pt}
  {بی اے \ \textbullet\ \ ایسوسی ایٹ ڈگری}\par
  \vspace{20pt}
  {\Huge\bfseries اسلامیات (اختیاری)}\par
  \vspace{6pt}
  {\large\bfseries \textenglish{Islamiat (Elective)}}\par
  \vspace{12pt}
  {\LARGE\bfseries متوقع سوالات اور ان کے جوابات}\par
  \vspace{12pt}
  {\large کورس کوڈ \textenglish{437}}\par
  \vspace{4pt}
  {سولہ سوالات \ \textbullet\ \ چار حقیقی پرچوں سے ماخوذ\par}
}

\vspace{1.0cm}

\begin{tcolorbox}[enhanced, colback=bmtint, colframe=bmblue, boxrule=0.6pt,
                  arc=3pt, left=12pt, right=12pt, top=10pt, bottom=10pt]
\textbf{\color{bmblue}سوالات کہاں سے لیے گئے ہیں؟}\ \
تمام سوالات اے آئی او یو کے \textbf{چار حقیقی گزشتہ پرچوں} سے لیے گئے ہیں:

\smallskip
\begin{itemize}\setlength{\itemsep}{0pt}\setlength{\topsep}{2pt}
  \item بی اے --- سمسٹر \textbf{بہار \textenglish{2012}} اور \textbf{بہار
        \textenglish{2013}}
  \item ایسوسی ایٹ ڈگری --- سمسٹر \textbf{خزاں \textenglish{2021}} اور
        \textbf{بہار \textenglish{2022}}
\end{itemize}

\smallskip
دو کے بجائے چار پرچے اس لیے لیے گئے کہ \textenglish{2013} اور
\textenglish{2021} کے درمیان یہ کورس بی اے سے ایسوسی ایٹ ڈگری میں منتقل ہوا۔
صرف پرانے دو پرچے دیکھنے سے موجودہ رجحان کا غلط اندازہ ہوتا۔ بتیس سوالات
یکجا کر کے \textbf{سولہ الگ موضوعات} بنے۔

\medskip
\textbf{\color{bmblue}صحتِ مواد۔}\ \
ہر آیت کا حوالہ سورت اور آیت کے نمبر سے دیا گیا ہے، اور ہر حدیث کے ساتھ اس
کی کتاب کا نام ہے۔ فقہی مسائل میں پہلے \textbf{فقہ حنفی} کا موقف لکھا گیا ہے
--- یہی پاکستان اور اے آئی او یو کا معمول ہے --- اور جہاں ائمہ کا اختلاف ہے،
اسے چھپایا نہیں بلکہ صاف بیان کر دیا گیا ہے۔

\medskip
\textbf{\color{bmblue}حیثیت۔}\ \
ذاتی مطالعے کے لیے۔ یہ اے آئی او یو کی مطبوعہ کتاب نہیں۔ اسے من و عن نقل کر کے
اسائنمنٹ میں جمع کرانا منع ہے۔
\end{tcolorbox}

\clearpage

%---------------------------------------------------------------------
\section*{پرچے کا خاکہ}
\markboth{پرچے کا خاکہ}{}

پرچے کی ساخت دونوں ادوار میں تقریباً ایک جیسی ہے، مگر \textbf{دو تبدیلیاں}
ایسی ہیں جو نظر انداز نہیں ہونی چاہئیں:

\medskip
\begin{center}
\begin{tabular}{@{}p{4.4cm} p{4.6cm} p{5.0cm}@{}}
\toprule
 & \textbf{بی اے (\textenglish{2012}، \textenglish{2013})}
 & \textbf{ایسوسی ایٹ ڈگری (\textenglish{2021}، \textenglish{2022})} \\
\midrule
کل نمبر            & \textenglish{100} & \textenglish{100} \\
کامیابی کے نمبر     & \textenglish{40}  & \textbf{\textenglish{50}} \\
وقت                & تین گھنٹے          & تین گھنٹے \\
سوالات دیے جاتے ہیں & آٹھ               & آٹھ \\
حل کرنے ہیں         & پانچ              & پانچ \\
پہلا سوال           & لازمی             & \textenglish{2021} میں لازمی، \textenglish{2022} میں نہیں \\
فی سوال نمبر        & \textenglish{20}  & \textenglish{20} \\
\bottomrule
\end{tabular}
\end{center}

\medskip
\begin{ghalti}[سب سے مہنگی غلط فہمی]
کامیابی کے نمبر اب \textbf{چالیس نہیں، پچاس} ہیں۔ جو طالب علم پرانے پرچوں کی
بنیاد پر ''چالیس کافی ہیں'' سمجھ کر تیاری کرتا ہے، وہ دس نمبر کے فرق سے فیل
ہوتا ہے۔ اسی طرح پہلا سوال \emph{کبھی لازمی ہوتا ہے اور کبھی نہیں} --- پرچہ
ہاتھ میں آتے ہی سب سے پہلے \textbf{نوٹ} کی سطر پڑھیے، فرض نہ کیجیے۔
\end{ghalti}

\medskip
\noindent
\textbf{\color{bmblue}کون سا موضوع کتنے پرچوں میں آیا}

\medskip
\begin{center}
\begin{tabular}{@{}p{6.8cm} p{2.4cm} p{5.0cm}@{}}
\toprule
\textbf{موضوع} & \textbf{کتنے پرچوں میں} & \textbf{اس کتابچے میں} \\
\midrule
آل عمران کی آیات کا ترجمہ و تشریح & \textbf{چاروں} & سوال \textenglish{2} \\
احادیث کا ترجمہ و تشریح            & \textbf{چاروں} & سوال \textenglish{7--10} \\
سورۃ آل عمران پر موضوعاتی نوٹ      & \textbf{چاروں} & سوال \textenglish{1}، \textenglish{3--6} \\
اسلامی تاریخ --- اموی و عباسی       & \textbf{چاروں} & سوال \textenglish{14--16} \\
حفظ اللسان --- زبان کی حفاظت        & تین            & سوال \textenglish{7} \\
نماز کے مسائل                       & تین            & سوال \textenglish{11--13} \\
صلہ رحمی --- البر والصلۃ            & تین            & سوال \textenglish{8} \\
مختصر اصطلاحات (لازمی سوال)         & دو             & حصہ اول \\
الحب فی اللہ                        & دو             & سوال \textenglish{9} \\
\bottomrule
\end{tabular}
\end{center}

\medskip
\noindent
پہلی چار سطریں \emph{ہر} پرچے میں آئی ہیں۔ صرف انہی چار موضوعات سے عموماً
پانچ میں سے تین یا چار سوال بن جاتے ہیں۔

\medskip
\begin{nukta}[اسلامیات کے جواب کا سانچہ]
اس مضمون میں ممتحن \emph{حوالہ} دیکھتا ہے۔ صرف مضمون لکھ دینا آدھے نمبر دیتا
ہے۔ ہر جواب میں:

\smallskip
\textbf{(۱)} تعریف یا تعارف،
\textbf{(۲)} کم از کم ایک \textbf{آیت} حوالے کے ساتھ،
\textbf{(۳)} کم از کم ایک \textbf{حدیث} کتاب کے نام کے ساتھ،
\textbf{(۴)} تفصیل \emph{عنوانات کے تحت}، اور
\textbf{(۵)} مختصر خلاصہ۔

\smallskip
عربی عبارت \textbf{خوش خط} اور اعراب کے ساتھ لکھیے --- بہار \textenglish{2012}
کے پرچے میں تو ''خوش خط لکھ کر'' الفاظ ہی سوال کا حصہ تھے۔ اگر پوری آیت یاد نہ
ہو تو جتنی یاد ہے اتنی صحیح لکھیے؛ غلط عربی لکھنے سے نہ لکھنا بہتر ہے۔
\end{nukta}
